\documentclass[aspectratio=169,10pt]{beamer}

\usepackage{iftex}
\ifPDFTeX
  \usepackage[T1]{fontenc}
  \usepackage[utf8]{inputenc}
  \usepackage{lmodern}
\fi
\usepackage[spanish,es-noshorthands]{babel}
\usepackage{booktabs}
\usepackage{array}
\usepackage{graphicx}
\usepackage{qrcode}
\usepackage{tikz}
\usetikzlibrary{arrows.meta,positioning,shapes.geometric,calc,fit,babel}
\usepackage{pgfplots}
\usepgfplotslibrary{statistics}

% Paleta común para todos los diagramas
\colorlet{cSensor}{blue!60!black}
\colorlet{cProc}{orange!80!black}
\colorlet{cAct}{green!45!black}
\colorlet{cApp}{violet!70!black}
\colorlet{cPower}{gray!70!black}

% Colores de los gráficos de datos (paleta categórica validada para daltonismo)
\definecolor{serie1}{HTML}{2A78D6}
\definecolor{serie2}{HTML}{EB6834}
\definecolor{serie3}{HTML}{1BAF7A}
\definecolor{tinta}{HTML}{52514E}

\tikzset{
  >={Stealth[length=2mm]},
  flecha/.style={->, thick, draw=black!75},
  caja/.style={rectangle, rounded corners=2pt, draw=#1, fill=#1!10, thick,
    minimum width=2.3cm, minimum height=0.85cm, align=center, inner sep=2pt},
  etiqueta/.style={font=\scriptsize, text=black!70, fill=white, inner sep=1pt},
}

% Estilo base de los gráficos. Las diapositivas cambian el tamaño con \renewcommand.
\newcommand{\anchografico}{9cm}
\newcommand{\altografico}{5.2cm}
\pgfplotsset{
  compat=1.18,
  grafico/.style={
    axis lines=left,
    axis line style={draw=black!45},
    tick style={draw=black!45},
    ymajorgrids, grid style={draw=black!12},
    tick label style={font=\scriptsize, text=tinta},
    label style={font=\footnotesize, text=tinta},
    legend style={draw=none, fill=none, font=\scriptsize, text=tinta},
    legend cell align=left,
    every axis plot/.append style={line width=0.8pt},
    /pgf/number format/use comma,
  },
}

% Mientras los archivos de figuras/datos/ tengan valores de ejemplo, cada
% gráfico lleva la marca "Datos de ejemplo". Al cargar las mediciones reales,
% cambiar a \datosdeejemplofalse en main.tex y en slides.tex.
\newif\ifdatosdeejemplo
\datosdeejemplotrue
\newcommand{\marcaejemplo}{%
  \ifdatosdeejemplo
    \node[anchor=north east, font=\scriptsize\bfseries, text=red!70!black,
      fill=white, draw=red!70!black, rounded corners=1pt, inner sep=2pt]
      at (rel axis cs:0.98,0.98) {Datos de ejemplo};
  \fi}

% Dato que falta medir o confirmar. Sale en rojo para que no pase desapercibido.
\newcommand{\pendiente}[1]{\textcolor{red!70!black}{[#1]}}

% Foto o captura: \imagen{archivo}{ancho}{alto}{qué debe mostrar}
% Si Imagenes/archivo existe, se inserta. Si no, se dibuja un recuadro con la
% descripción, así el documento compila mientras faltan las fotos.
\newcommand{\imagen}[4]{%
  \IfFileExists{Imagenes/#1}{%
    \includegraphics[width=#2, height=#3, keepaspectratio]{Imagenes/#1}%
  }{%
    \begin{tikzpicture}
      \node[draw=black!35, dashed, fill=black!4, rounded corners=2pt,
        minimum width=#2, minimum height=#3, text width=0.85\dimexpr#2\relax,
        align=center, font=\scriptsize, text=black!60] {%
        \textbf{Imagenes/#1}\\[1mm]#4};
    \end{tikzpicture}%
  }}

% Gráficos más chicos que en el informe
\renewcommand{\anchografico}{7.2cm}
\renewcommand{\altografico}{4.6cm}

% Enlace del video publicitario. Si se deja vacío, la lámina muestra solo la imagen.
\def\videourl{}

% ---------------------------------------------------------------
% Tema académico sobrio
\usetheme{Boadilla}
\definecolor{azul}{HTML}{1F3A5F}
\definecolor{gris}{HTML}{5F6B7A}
\usecolortheme[named=azul]{structure}
\setbeamercolor{block title}{bg=azul!12, fg=azul}
\setbeamercolor{block body}{bg=azul!4}
\setbeamercolor{alerted text}{fg=cProc}
\setbeamertemplate{blocks}[rounded][shadow=false]
\setbeamertemplate{navigation symbols}{}
\setbeamertemplate{itemize items}[circle]
\setbeamertemplate{bibliography item}[text]
\setbeamerfont{frametitle}{series=\bfseries}

% Estilos TikZ propios de las diapositivas (fuera de los frames por el #)
\tikzset{
  capa/.style={rectangle, minimum width=2.4cm, minimum height=0.9cm,
    fill=#1, text=white, font=\small\bfseries},
  nuestro/.style={rectangle, draw=#1, fill=#1!8, thick, minimum height=0.9cm,
    align=left, font=\small, text width=3.2cm},
  n/.style={caja=#1, minimum width=1.8cm, minimum height=1cm},
  antes/.style={caja=cPower, minimum width=3.1cm, minimum height=0.75cm, font=\small},
  ahora/.style={caja=cAct, minimum width=3.1cm, minimum height=0.75cm, font=\small},
  paso/.style={circle, fill=azul, text=white, font=\small\bfseries, minimum size=6mm, inner sep=0pt},
}

% Foto de referencia con leyenda: \foto{archivo}{alto}{título}{detalle}
\newcommand{\foto}[4]{%
  \begin{minipage}[t]{\linewidth}\centering
    \fcolorbox{gris!40}{white}{\begin{minipage}[c][#2][c]{0.92\linewidth}\centering
      \includegraphics[height=0.94\dimexpr#2\relax, width=0.94\linewidth, keepaspectratio]{figuras/fotos/#1}
    \end{minipage}}\\[1mm]
    {\footnotesize\bfseries #3}\\[-0.5mm]
    {\scriptsize\color{gris} #4}
  \end{minipage}}
% Foto propia con leyenda (placeholder si aún no está en Imagenes/)
\newcommand{\fotopropia}[5]{%
  \begin{minipage}[t]{\linewidth}\centering
    \imagen{#1}{0.95\linewidth}{#2}{#5}\\[1mm]
    {\footnotesize\bfseries #3}\\[-0.5mm]
    {\scriptsize\color{gris} #4}
  \end{minipage}}
\setlength{\fboxsep}{0pt}
\setlength{\fboxrule}{0.3pt}

% Crédito de imágenes, sobre el pie de página
\newcommand{\credito}[1]{%
  \begin{tikzpicture}[remember picture, overlay]
    \node[anchor=south west, font=\tiny, text=gris, inner sep=0pt, text width=15cm]
      at ([xshift=4mm, yshift=6.5mm]current page.south west) {#1};
  \end{tikzpicture}}

\title[Seguridad residencial IoT]{Sistema inteligente de seguridad residencial IoT\\
  con monitoreo inalámbrico y alertas en tiempo real}
\subtitle{Presentación final (P2)}
\author[Salcedo, Quispe, Iribar, Salazar]{Ariann Abigail Salcedo Quiroz \and
  Diego Randolp Quispe Apaza \\ Federico Iribar Casanova \and
  Sofia del Pilar Salazar Mendoza}
\institute[UTEC]{Internet de las Cosas\\Universidad de Ingeniería y Tecnología}
\date[Nov. 2026]{Noviembre de 2026}

\begin{document}

\begin{frame}
  \titlepage
\end{frame}

% ===============================================================
\section{Introducción}

\begin{frame}{Introducción}
  \begin{itemize}
    \item Una vivienda suele protegerse con piezas que no se hablan: cerradura, alarma y, con suerte, una cámara.
    \item Cada una trabaja sola. La casa tarda en notar un riesgo y casi nunca reacciona por sí misma.
    \item Hoy un ESP32 de pocos dólares puede leer sensores, decidir y avisar en tiempo real \cite{alfuqaha2015survey, oktafiani2025esp32}.
  \end{itemize}
  \vspace{4mm}
  \begin{block}{Nuestra propuesta en cinco etapas}
    \centering
    \begin{tikzpicture}[font=\small, node distance=4mm]
      \node[caja=cSensor, minimum width=2.3cm] (a) {Detectar\\{\scriptsize 3 sensores}};
      \node[caja=cProc, minimum width=2.3cm, right=of a] (b) {Decidir\\{\scriptsize ESP32, 10\,s}};
      \node[caja=cAct, minimum width=2.3cm, right=of b] (c) {Responder\\{\scriptsize buzzer y pestillo}};
      \node[caja=cApp, minimum width=2.3cm, right=of c] (d) {Avisar\\{\scriptsize MQTT a AWS}};
      \node[caja=cApp, minimum width=2.3cm, right=of d] (e) {Medir\\{\scriptsize 6 métricas}};
      \foreach \u/\v in {a/b, b/c, c/d, d/e} \draw[flecha] (\u) -- (\v);
    \end{tikzpicture}
  \end{block}
\end{frame}

% ===============================================================
\section{Problemática}

\begin{frame}{Problemática}
  Una intrusión puede empezar abriendo una puerta, con movimiento en una zona vacía o con golpes contra una estructura \cite{arinde2024multisensor}.
  \vspace{3mm}
  \begin{columns}[T]
    \column{0.32\textwidth}
    \begin{block}{Desconexión}
      Nadie relaciona una apertura con un golpe de hace 5 segundos. Se pierden los primeros segundos.
    \end{block}
    \column{0.32\textwidth}
    \begin{block}{Costo}
      Lo comercial es caro. Lo barato trae un solo sensor y rara vez avisa a distancia \cite{oktafiani2025esp32}.
    \end{block}
    \column{0.32\textwidth}
    \begin{block}{Sin métricas}
      El dueño no sabe cuántas alertas fueron falsas ni cuánto tardó un aviso \cite{pugh2015fusion}.
    \end{block}
  \end{columns}
  \vspace{3mm}
  \begin{alertblock}{Problema}
    Falta un sistema integrado y barato que correlacione varios sensores, responda sin depender de Internet y mida su propio desempeño.
  \end{alertblock}
\end{frame}

% ===============================================================
\section{Marco teórico}

\begin{frame}{Marco teórico}
  \begin{columns}[T]
    \column{0.47\textwidth}
    \textbf{Arquitectura IoT por capas} \cite{alfuqaha2015survey, ray2018survey}\par\vspace{3mm}
    \begin{tikzpicture}[node distance=2mm]
      \node[capa=cApp] (ap) {Aplicación};
      \node[capa=cProc, below=of ap] (rd) {Red};
      \node[capa=cSensor, below=of rd] (pe) {Percepción};
      \node[nuestro=cApp, right=3mm of ap] {AWS y app móvil};
      \node[nuestro=cProc, right=3mm of rd] {Wi-Fi + MQTT/TLS};
      \node[nuestro=cSensor, right=3mm of pe] {Sensores y actuadores};
    \end{tikzpicture}\par\vspace{4mm}
    \small
    \begin{itemize}
      \item MQTT publica y suscribe con 2 bytes de cabecera \cite{light2017mosquitto}.
      \item Lo urgente se decide en el borde \cite{shi2016edge}.
    \end{itemize}
    \column{0.5\textwidth}
    \textbf{Sensores}\par\vspace{2mm}
    \begin{columns}[T]
      \column{0.32\linewidth}\foto{pir.jpg}{2.4cm}{PIR HC-SR501}{Movimiento, 3 a 7\,m}
      \column{0.32\linewidth}\foto{reed.jpg}{2.4cm}{Magnético MC-38}{Apertura de la puerta}
      \column{0.32\linewidth}\foto{sw420.jpg}{2.4cm}{Vibración SW-420}{Golpes}
    \end{columns}
    \vspace{5mm}
    \small
    \begin{itemize}
      \item AWS IoT Core es un broker MQTT administrado que autentica con certificados \cite{aws2026iot}.
    \end{itemize}
  \end{columns}
  \credito{Fotos referenciales de Wikimedia Commons: PIR por Lethalattraction (CC BY-SA 4.0), reed switch por Satishbhandare (dominio público), SW-420 por Suyash Dwivedi (CC BY-SA 4.0).}
\end{frame}

% ===============================================================
\section{Estado del arte}

\begin{frame}{Estado del arte (2020 a 2026)}
  \footnotesize
  \renewcommand{\arraystretch}{1.15}
  \begin{tabular}{@{}l>{\raggedright\arraybackslash}p{7.4cm}>{\raggedright\arraybackslash}p{4.7cm}@{}}
    \toprule
    \textbf{Integrante} & \textbf{Trabajo} & \textbf{Qué tomamos} \\
    \midrule
    Ariann & Nannung et al. \cite{nannung2025iot}. ESP32-CAM, rostro al 95,31\,\%, relé y cerradura & Relé como respuesta física \\
     & Oktafiani y Widhiantoro \cite{oktafiani2025esp32}. PIR en 7 de 10 sistemas ESP32 & PIR como sensor principal \\
    \midrule
    Diego & Aprilia y Purwanto \cite{aprilia2025boarding}. PIR + buzzer, 100\,\% hasta 3\,m, aviso en 8 a 9\,s & Ubicación del PIR, medir el aviso \\
     & Jama et al. \cite{jama2026edge}. ESP32 decide en el borde y sigue con la red caída & Alarma local sin nube \\
    \midrule
    Federico & Mishra y Kertesz \cite{mishra2020mqtt}. MQTT es el protocolo que más crece & Elección de MQTT \\
     & Bender et al. \cite{bender2021mqtt}. Seis brokers comparados en latencia y recursos & Cliente MQTT estándar, medir la latencia propia \\
    \midrule
    Sofia & Zhuang et al. \cite{zhuang2021intrusion}. Intrusos detectados por CSI de Wi-Fi en casas reales & Trabajo futuro, RSSI como métrica \\
     & Arinde e Idowu \cite{arinde2024multisensor}. PIR + magnético + llamada GSM & La misma base de sensores \\
    \bottomrule
  \end{tabular}
\end{frame}

\begin{frame}{Del estado del arte al diseño}
  \centering
  \resizebox{\linewidth}{!}{% Qué tomamos del estado del arte en cada etapa del prototipo, y qué agregamos.
\begin{tikzpicture}[font=\footnotesize,
  etapa/.style={caja=#1, minimum width=2.2cm, minimum height=0.8cm, font=\footnotesize\bfseries},
  ref/.style={rectangle, rounded corners=1.5pt, draw=black!30, fill=white,
    font=\scriptsize, minimum width=2.2cm, minimum height=0.5cm},
  nuevo/.style={rectangle, rounded corners=1.5pt, draw=#1, fill=#1!15, thick,
    font=\scriptsize, minimum width=2.2cm, minimum height=0.55cm, align=center},
]
  \node[etapa=cSensor] (s1) at (0,0)    {Detectar};
  \node[etapa=cProc]   (s2) at (2.6,0)  {Decidir};
  \node[etapa=cAct]    (s3) at (5.2,0)  {Responder};
  \node[etapa=cApp]    (s4) at (7.8,0)  {Avisar};
  \node[etapa=cApp]    (s5) at (10.4,0) {Medir};
  \foreach \a/\b in {s1/s2, s2/s3, s3/s4, s4/s5} \draw[flecha] (\a) -- (\b);

  % Trabajos previos que inspiran cada etapa
  \node[ref] at (0,1.05)    {PIR + magnético \cite{arinde2024multisensor}};
  \node[ref] at (0,1.65)    {PIR en 7 de 10 \cite{oktafiani2025esp32}};
  \node[ref] at (2.6,1.05)  {Decidir en el borde \cite{jama2026edge}};
  \node[ref] at (5.2,1.05)  {Cerradura por orden \cite{nannung2025iot}};
  \node[ref] at (5.2,1.65)  {Buzzer local \cite{aprilia2025boarding}};
  \node[ref] at (7.8,1.05)  {MQTT \cite{mishra2020mqtt}};
  \node[ref] at (7.8,1.65)  {Brokers compatibles \cite{bender2021mqtt}};
  \node[ref] at (10.4,1.05) {Tiempo de aviso \cite{aprilia2025boarding}};
  \node[ref] at (10.4,1.65) {Señal Wi-Fi \cite{zhuang2021intrusion}};
  \node[font=\scriptsize\itshape, text=black!55, anchor=east] at (-1.25,1.35) {Estado del arte};

  % Lo que agrega este proyecto
  \node[nuevo=cProc] at (2.6,-1.0)  {Correlación\\de 2 sensores};
  \node[nuevo=cAct]  at (5.2,-1.0)  {Pestillo con PIN que\\se traba al cerrar};
  \node[nuevo=cApp]  at (10.4,-1.0) {6 métricas\\propias};
  \node[font=\scriptsize\itshape, text=black!55, anchor=east] at (-1.25,-1.0) {Aporte propio};
\end{tikzpicture}
}
  \vspace{5mm}
  \begin{alertblock}{Lo que agregamos}
    \centering Dos sensores deben coincidir, el pestillo se suelta con una orden y se traba al cerrar, y el sistema se mide a sí mismo.
  \end{alertblock}
\end{frame}

% ===============================================================
\section{Metodología}

\begin{frame}{Qué cambió desde el perfil}
  \begin{columns}[c]
    \column{0.52\textwidth}
    \centering
    \begin{tikzpicture}[font=\small]
      \foreach \y/\a/\b in {0/Sirena de 12\,V/Buzzer SFM-20B, -0.8/Cerradura solenoide/Servo SG90 + pestillo, -1.6/Fuente de 12\,V + LM2596/Cable USB de 5\,V, -2.4/Placa perforada/Protoboard, -3.2/Carcasa impresa en 3D/Casa de cartón reciclado, -4.0/Mosquitto en EC2/AWS IoT Core + Lambda, -4.8/App web/App móvil} {
        \node[antes] (x) at (0,\y) {\a};
        \node[ahora] (y) at (4.3,\y) {\b};
        \draw[flecha] (x) -- (y);
      }
      \node[font=\scriptsize\bfseries, text=cPower] at (0,0.75) {Perfil P1};
      \node[font=\scriptsize\bfseries, text=cAct] at (4.3,0.75) {Versión final};
    \end{tikzpicture}
    \column{0.46\textwidth}
    \centering
    \resizebox{\linewidth}{!}{% Potencia estimada en el peor caso (W).
% Perfil P1, a 12 V (hoja "Presupuesto de corriente"): electrónica 0,36 A = 4,3 W;
% sirena + cerradura 1,90 A = 22,8 W.
% Diseño final, a 5 V por USB: ESP32 0,50 A + PIR 0,065 A + SW-420 0,015 A + LED 0,01 A
% = 0,59 A = 2,95 W; buzzer 0,03 A + servo atascado 0,70 A = 0,73 A = 3,65 W;
% MQ-2 0,15 A = 0,75 W.
\begin{tikzpicture}
\begin{axis}[grafico, ybar stacked, bar width=1.1cm,
  width=\anchografico, height=\altografico,
  symbolic x coords={P1,Final,Mq},
  xtick=data, xticklabels={{Perfil P1\\27,1\,W},{Final\\6,6\,W},{Final + MQ-2\\7,4\,W}},
  xticklabel style={align=center},
  enlarge x limits=0.25, ymin=0, ymax=30, ytick={0,5,...,30},
  ylabel={Potencia (W)},
  legend style={at={(0.5,1.03)}, anchor=south, legend columns=3},
  reverse legend=false]
  \addplot[fill=serie1, draw=white] coordinates {(P1,4.3) (Final,2.95) (Mq,2.95)};
  \addplot[fill=serie2, draw=white] coordinates {(P1,22.8) (Final,3.65) (Mq,3.65)};
  \addplot[fill=serie3, draw=white] coordinates {(P1,0) (Final,0) (Mq,0.75)};
  \legend{Electrónica, Actuadores, MQ-2}
  \draw[dashed, black!60, thick] ({rel axis cs:0,0} |- {axis cs:P1,24}) -- ({rel axis cs:1,0} |- {axis cs:P1,24})
    node[pos=0.98, above left, font=\scriptsize, text=tinta] {Fuente de 12\,V y 2\,A};
  \draw[dashed, black!60, thick] ({rel axis cs:0,0} |- {axis cs:P1,10}) -- ({rel axis cs:1,0} |- {axis cs:P1,10})
    node[pos=0.98, above left, font=\scriptsize, text=tinta] {Cargador USB de 5\,V y 2\,A};
\end{axis}
\end{tikzpicture}
}\par
    {\footnotesize Peor caso estimado: de 27\,W a unos 7\,W.}
  \end{columns}
\end{frame}

\begin{frame}{Componentes}
  \begin{columns}[T]
    \column{0.2\textwidth}\foto{esp32.jpg}{2.2cm}{ESP32}{Decide y publica}
    \column{0.36\textwidth}\fotopropia{buzzer_sg90.jpeg}{2.2cm}{Buzzer y servo SG90}{Alarma y pestillo}{Foto del buzzer SFM-20B y del servo SG90}
    \column{0.2\textwidth}\foto{esp32cam.jpg}{2.2cm}{ESP32-CAM}{Opcional}
  \end{columns}
  \vspace{4mm}
  \begin{columns}[T]
    \column{0.48\textwidth}
    \begin{block}{Núcleo}
      \small
      \begin{itemize}
        \item Sensores PIR, MC-38 y SW-420, y MQ-2 para humo
        \item Buzzer, servo y dos LED de estado
        \item Wi-Fi + MQTT/TLS hacia AWS IoT Core
        \item App móvil con apertura por PIN
      \end{itemize}
    \end{block}
    \column{0.48\textwidth}
    \begin{block}{ESP32 DevKit V1}
      \small
      \begin{itemize}
        \item 240\,MHz, 520\,KB de SRAM, Wi-Fi de 2,4\,GHz
        \item Hora NTP para medir la latencia
        \item S/~40, barato como se recomendó
      \end{itemize}
    \end{block}
  \end{columns}
  \credito{Fotos referenciales de Wikimedia Commons: ESP32 por Edwiyanto y ESP32-CAM por Nowforever (CC BY-SA 4.0).}
\end{frame}

\begin{frame}{Diagrama de bloques}
  \begin{columns}[c]
    \column{0.66\textwidth}
    \centering\resizebox{\linewidth}{!}{\begin{tikzpicture}[font=\footnotesize]
  % Percepción
  \node[caja=cSensor] (pir) at (-2.7,0)  {PIR\\{\scriptsize movimiento}};
  \node[caja=cSensor] (mag) at (0,0)     {MC-38\\{\scriptsize apertura}};
  \node[caja=cSensor] (vib) at (2.7,0)   {SW-420\\{\scriptsize golpes}};
  % Procesamiento
  \node[caja=cProc, minimum height=1cm] (esp) at (0,-1.9) {\textbf{ESP32}\\{\scriptsize regla de 10\,s}};
  % Actuación
  \node[caja=cAct] (relay) at (-1.6,-3.7) {Relé};
  \node[caja=cAct] (buz)   at (-1.6,-5.3) {Buzzer\\{\scriptsize alarma}};
  \node[caja=cAct] (servo) at (1.6,-3.7)  {Servo SG90\\{\scriptsize pestillo}};
  % Energía
  \node[caja=cPower, minimum width=1.9cm] (lm)  at (-4.9,-1.9) {LM2596\\{\scriptsize 12 a 5\,V}};
  \node[caja=cPower, minimum width=1.9cm] (psu) at (-4.9,-3.7) {Fuente\\{\scriptsize 12\,V, 2\,A}};
  % Nube
  \node[caja=cApp, minimum width=2.1cm] (brk) at (5.4,-1.9) {IoT Core\\{\scriptsize broker MQTT}};
  \node[caja=cApp, minimum width=2.1cm] (db)  at (5.4,-3.7) {DynamoDB};
  \node[caja=cApp, minimum width=2.1cm] (lam) at (8.4,-1.9) {Lambda};
  \node[caja=cApp, minimum width=2.1cm] (api) at (8.4,-3.7) {API Gateway};
  \node[caja=cApp, minimum width=2.1cm] (web) at (8.4,-5.3) {App móvil};
  \node[draw=cApp, dashed, rounded corners=3pt, inner sep=5pt, fit=(brk)(db)(lam)(api),
    label={[font=\scriptsize, text=cApp]above:AWS}] {};

  % Sensores al ESP32
  \draw[flecha] (pir.south) -- ++(0,-0.45) -| ([xshift=-7mm]esp.north);
  \draw[flecha] (mag.south) -- (esp.north);
  \draw[flecha] (vib.south) -- ++(0,-0.45) -| ([xshift=7mm]esp.north);
  % ESP32 a actuadores
  \draw[flecha] ([xshift=-5mm]esp.south) |- ++(0,-0.4) -| (relay.north);
  \draw[flecha] ([xshift=5mm]esp.south) |- ++(0,-0.4) -| (servo.north);
  \draw[flecha] (relay) -- (buz);
  % Energía
  \draw[flecha, draw=cPower] (psu) -- node[etiqueta, left] {12\,V} (lm);
  \draw[flecha, draw=cPower] (lm) -- node[etiqueta, above] {5\,V} (esp);
  \draw[flecha, draw=cPower] (psu) -- node[etiqueta, above] {12\,V} (relay);
  % Red y nube
  \draw[flecha, <->, dashed, draw=cApp] (esp) -- node[etiqueta, above, text=cApp] {Wi-Fi} node[etiqueta, below, text=cApp] {MQTT/TLS} (brk);
  \draw[flecha, draw=cApp] (brk) -- node[etiqueta, left, text=cApp] {reglas} (db);
  \draw[flecha, draw=cApp] (lam) -- node[etiqueta, above, text=cApp] {abrir} (brk);
  \draw[flecha, <->, draw=cApp] (lam) -- (db);
  \draw[flecha, draw=cApp] (api) -- (lam);
  \draw[flecha, <->, draw=cApp] (web) -- node[etiqueta, right, text=cApp] {HTTPS} (api);
\end{tikzpicture}
}
    \column{0.32\textwidth}
    \small
    \begin{itemize}
      \item {\color{cSensor}\textbf{Azul}} sensores
      \item {\color{cProc}\textbf{Naranja}} decisión local
      \item {\color{cAct}\textbf{Verde}} actuación
      \item {\color{cPower}\textbf{Gris}} energía
      \item {\color{cApp}\textbf{Morado}} nube e interfaz
    \end{itemize}
  \end{columns}
\end{frame}

% ===============================================================
\section{Implementación}

\begin{frame}{Etapa 1 · Percepción e interconexiones}
  \begin{columns}[c]
    \column{0.62\textwidth}
    \centering
    \resizebox{\linewidth}{!}{% Interconexiones del nodo: qué pin del ESP32 va a cada módulo.
\begin{tikzpicture}[font=\footnotesize,
  pin/.style={font=\scriptsize\ttfamily, text=black!80},
  nota/.style={font=\scriptsize, text=black!55},
  senal/.style={thick, draw=black!75},
  v12/.style={thick, draw=red!65!black},
]
  % ESP32 con sus pines
  \node[caja=cProc, minimum width=2.6cm, minimum height=4.4cm] (esp) at (0,0) {};
  \node[font=\footnotesize\bfseries, text=cProc] at (0,1.85) {ESP32};
  \node[pin, anchor=west] (p27) at (-1.25,1.0)  {GPIO27};
  \node[pin, anchor=west] (p26) at (-1.25,0.0)  {GPIO26};
  \node[pin, anchor=west] (p25) at (-1.25,-1.0) {GPIO25};
  \node[pin, anchor=east] (p23) at (1.25,1.0)   {GPIO23};
  \node[pin, anchor=east] (p18) at (1.25,-0.4)  {GPIO18};
  \node[pin, anchor=east] (pvin) at (1.25,-1.9) {VIN};
  \node[pin, anchor=west] (p33) at (-1.25,-1.8) {3V3};

  % Sensores a la izquierda
  \node[caja=cSensor, minimum width=1.9cm] (pir) at (-4.3,1.0)  {PIR};
  \node[caja=cSensor, minimum width=1.9cm] (mag) at (-4.3,0.0)  {MC-38};
  \node[caja=cSensor, minimum width=1.9cm] (vib) at (-4.3,-1.0) {SW-420};
  \draw[senal] (pir.east) -- (esp.west |- pir.east);
  \draw[senal] (mag.east) -- (esp.west |- mag.east);
  \draw[senal] (vib.east) -- (esp.west |- vib.east);
  \node[nota, left=1pt of pir] {5\,V};
  \node[nota, left=1pt of mag] {GND};
  \node[nota, left=1pt of vib] {3,3\,V};

  % Actuadores a la derecha
  \node[caja=cAct, minimum width=1.9cm] (rel) at (4.1,1.0)  {Relé CH1\\{\scriptsize jumper en H}};
  \node[caja=cAct, minimum width=1.9cm] (srv) at (4.1,-0.4) {Servo SG90};
  \node[caja=cAct, minimum width=1.9cm] (buz) at (7.3,1.0)  {Buzzer};
  \draw[senal] (esp.east |- rel.west) -- node[etiqueta, above] {IN1} (rel.west);
  \draw[senal] (esp.east |- srv.west) -- node[etiqueta, above] {PWM} (srv.west);
  \draw[v12] (rel.east) -- node[etiqueta, above] {NO} (buz.west);
  \draw[v12] (rel.north) -- ++(0,0.5) node[above, font=\scriptsize, text=red!65!black] {12\,V a COM};
  \node[nota, below=1pt of srv] {+ 1000\,\textmu F};

  % Alimentación de 5 V
  \node[caja=cPower, minimum width=1.9cm] (lm) at (4.1,-1.9) {LM2596\\{\scriptsize 5,0\,V}};
  \draw[thick, draw=cPower] (lm.west) -- (esp.east |- lm.west);
  \node[nota, anchor=north] at (0,-2.45) {Todas las tierras unidas};
\end{tikzpicture}
}
    \column{0.36\textwidth}
    \small
    \begin{itemize}
      \item SW-420 a 3,3\,V y todas las tierras unidas
      \item MQ-2 con divisor de 10\,k$\Omega$ y 20\,k$\Omega$ en el ADC1
      \item Buzzer directo al GPIO 32
      \item LED rojo y verde con 220\,$\Omega$
    \end{itemize}
  \end{columns}
\end{frame}

\begin{frame}{Etapa 2 · Decisión en el ESP32}
  \begin{columns}[c]
    \column{0.56\textwidth}
    \centering
    \resizebox{\linewidth}{!}{% Dos casos sobre una línea de tiempo: la ventana de 10 s decide el nivel.
% 1 cm = 2 s
\begin{tikzpicture}[font=\footnotesize, x=0.5cm,
  evento/.style={circle, fill=#1, draw=white, line width=0.8pt, minimum size=7pt, inner sep=0pt},
  nombre/.style={font=\scriptsize, text=black!75},
  resultado/.style={rounded corners=2pt, fill=#1!12, draw=#1, thick, font=\scriptsize\bfseries,
    text=#1, minimum width=1.7cm, minimum height=0.6cm},
]
  % Caso 1: golpe y luego apertura, 5 s de diferencia
  \node[anchor=east, font=\footnotesize\bfseries] at (-0.6,0) {Caso 1};
  \fill[cAct!12] (2,-0.32) rectangle (12,0.32);
  \draw[black!40] (0,0) -- (22,0);
  \node[evento=cSensor] (a1) at (2,0) {};
  \node[evento=cSensor] (a2) at (7,0) {};
  \node[nombre, above=2pt of a1] {SW-420};
  \node[nombre, above=2pt of a2] {MC-38};
  \draw[|<->|, cAct, thick] (2,-0.55) -- node[below, font=\scriptsize] {ventana de 10\,s} (12,-0.55);
  \node[resultado=cAct, anchor=west] at (23,0) {Crítico};

  % Caso 2: movimiento y apertura separados por 13 s
  \begin{scope}[yshift=-1.9cm]
    \node[anchor=east, font=\footnotesize\bfseries] at (-0.6,0) {Caso 2};
    \fill[cProc!10] (3,-0.32) rectangle (13,0.32);
    \draw[black!40] (0,0) -- (22,0);
    \node[evento=cSensor] (b1) at (3,0) {};
    \node[evento=cSensor] (b2) at (16,0) {};
    \node[nombre, above=2pt of b1] {PIR};
    \node[nombre, above=2pt of b2] {MC-38};
    \draw[|<->|, cProc, thick] (3,-0.55) -- node[below, font=\scriptsize] {10\,s} (13,-0.55);
    \node[resultado=cProc, anchor=west] at (23,0) {2 leves};
  \end{scope}

  % Eje de tiempo
  \begin{scope}[yshift=-3.3cm]
    \draw[->, black!60] (0,0) -- (22,0) node[right, font=\scriptsize, text=black!60] {$t$ (s)};
    \foreach \t in {0,5,10,15,20}
      \draw[black!60] (\t,0.08) -- (\t,-0.08) node[below, font=\scriptsize, text=black!60] {\t};
  \end{scope}
\end{tikzpicture}
}\par\vspace{5mm}
    \small
    \begin{itemize}
      \item Un sensor solo es una \textbf{alerta leve}
      \item Dos sensores distintos en 10\,s son un \textbf{evento crítico}
      \item La decisión se toma antes de publicar, sin esperar a la red
      \item Cada 30\,s sale un latido con el RSSI
    \end{itemize}
    \column{0.42\textwidth}
    \centering
    \resizebox{!}{0.8\textheight}{% Flujo del programa del ESP32. Arriba, el ciclo de lectura (núcleo 1); las cajas
% punteadas corren en paralelo (la tarea de red en el núcleo 0).
\begin{tikzpicture}[
  font=\footnotesize,
  inicio/.style={rectangle, rounded corners=3.5mm, draw=black!70, fill=black!8, thick,
    minimum width=2cm, minimum height=0.7cm},
  proceso/.style={caja=cSensor, minimum width=2.8cm, minimum height=0.8cm},
  accion/.style={caja=cAct, minimum width=2.8cm, minimum height=0.8cm},
  aviso/.style={caja=cProc, minimum width=2.4cm, minimum height=0.8cm},
  red/.style={caja=cApp, minimum width=2.8cm, minimum height=0.8cm},
  decision/.style={diamond, aspect=2, draw=cProc, fill=cProc!10, thick,
    align=center, inner sep=1pt},
  si/.style={font=\scriptsize, text=black!70},
]
  \node[inicio]   (start) at (0,0)     {Inicio};
  \node[proceso]  (conf)  at (0,-1.2)  {Bloquear la puerta,\\crear colas y red};
  \node[proceso]  (leer)  at (0,-2.6)  {Leer sensores};
  \node[decision] (d1)    at (0,-4.1)  {¿Sensor activo\\y armado?};
  \node[decision] (d2)    at (0,-5.9)  {¿Otro sensor\\en 10\,s?};
  \node[accion]   (crit)  at (0,-7.6)  {Buzzer 30\,s y\\bloquear puerta};
  \node[red]      (pub)   at (0,-9.1)  {Poner el evento\\en la cola};
  \node[aviso]    (leve)  at (3.6,-5.9) {Alerta leve};

  % Procesos paralelos
  \node[red, dashed]    (lat) at (3.8,0.3)  {Red (núcleo 0):\\publica la cola y\\un latido cada 30\,s};
  \node[accion, dashed] (cmd) at (3.8,-1.35) {Orden de abrir:\\desbloquear};
  \node[accion, dashed] (gas) at (3.8,-2.75) {Humo en el MQ-2:\\buzzer y desbloquear};
  \node[accion, dashed] (ab)  at (3.8,-4.15) {Puerta abierta 1\,min:\\buzzer hasta cerrar};

  \draw[flecha] (start) -- (conf);
  \draw[flecha] (conf) -- (leer);
  \draw[flecha] (leer) -- (d1);
  \draw[flecha] (d1) -- node[si, right] {Sí} (d2);
  \draw[flecha] (d2) -- node[si, right] {Sí} (crit);
  \draw[flecha] (crit) -- (pub);
  % Sin activación: vuelve a leer
  \draw[flecha] (d1.west) -- node[si, above] {No} ++(-0.55,0) |- (leer.west);
  % Un solo sensor: alerta leve y se publica
  \draw[flecha] (d2.east) -- node[si, above] {No} (leve.west);
  \draw[flecha] (leve.south) |- (pub.east);
  % Después de poner el evento en la cola, vuelve a leer
  \draw[flecha, draw=black!45] (pub.south) -- ++(0,-0.4) -| (-2.9,-2.6) -- (leer.west);
\end{tikzpicture}
}
  \end{columns}
\end{frame}

\begin{frame}{Etapa 3 · Actuación y energía}
  \begin{columns}[c]
    \column{0.55\textwidth}
    \centering
    \resizebox{\linewidth}{!}{% Árbol de alimentación: del cable USB a cada módulo.
\begin{tikzpicture}[font=\footnotesize,
  v5/.style={flecha, draw=cPower},
  v3/.style={flecha, draw=cProc},
  riel/.style={font=\scriptsize\bfseries},
]
  \node[caja=cPower, minimum width=1.9cm] (usb) at (0,0) {Cable USB\\{\scriptsize 5\,V}};

  \node[caja=cProc,   minimum width=1.9cm] (esp) at (3.6,1.6)   {ESP32};
  \node[caja=cSensor, minimum width=1.9cm] (pir) at (3.6,0.55)  {PIR};
  \node[caja=cSensor, minimum width=1.9cm] (gas) at (3.6,-0.5)  {MQ-2};
  \node[caja=cAct,    minimum width=1.9cm] (srv) at (3.6,-1.55) {Servo};
  \node[caja=cSensor, minimum width=1.9cm] (vib) at (7.0,1.6)   {SW-420};

  \draw[v5, -] (usb.east) -- ++(0.6,0) coordinate (j5);
  \draw[v5] (j5) |- (esp.west);
  \draw[v5] (j5) |- (pir.west);
  \draw[v5] (j5) |- (gas.west);
  \draw[v5] (j5) |- (srv.west);
  \draw[v3] (esp) -- (vib);

  \node[riel, text=cPower, anchor=south east] at (2.55,-1.5) {5\,V};
  \node[riel, text=cProc, above] at (5.3,1.6) {3,3\,V};
  \node[font=\scriptsize, text=black!55, align=center] at (5.3,-2.5)
    {El buzzer y los LED se alimentan\\desde sus pines GPIO.};
\end{tikzpicture}
}
    \column{0.42\textwidth}
    \small
    \begin{itemize}
      \item Todo funciona con los 5\,V del cable USB, hoy desde una laptop
      \item La puerta está trabada por defecto y se suelta con la orden de la app
      \item El MC-38 ve que se abrió y cerró, y el servo la traba otra vez
      \item Más de 1\,min abierta: buzzer hasta que se cierre
      \item Humo en el MQ-2: buzzer y la puerta se destraba
    \end{itemize}
  \end{columns}
\end{frame}

\begin{frame}{Etapa 4 · Nube en AWS, sin servidor propio}
  \begin{columns}[c]
    \column{0.55\textwidth}
    \centering
    \resizebox{\linewidth}{!}{%
    \begin{tikzpicture}[font=\small, node distance=8mm and 14mm]
      \node[n=cProc] (esp) {ESP32};
      \node[n=cApp, right=18mm of esp] (iot) {IoT Core};
      \node[n=cApp, below=of iot] (db) {DynamoDB};
      \node[n=cApp, right=of iot] (lam) {Lambda};
      \node[n=cApp, below=of lam] (api) {API Gateway};
      \node[n=cApp, below=of api] (app) {App móvil};
      \draw[flecha, <->, dashed, draw=cApp] (esp) -- node[etiqueta, above] {MQTT/TLS} (iot);
      \draw[flecha, draw=cApp] (iot) -- node[etiqueta, left] {reglas} (db);
      \draw[flecha, draw=cApp] (lam) -- node[etiqueta, above] {abrir} (iot);
      \draw[flecha, <->, draw=cApp] (lam) -- (db);
      \draw[flecha, draw=cApp] (api) -- (lam);
      \draw[flecha, <->, draw=cApp] (app) -- (api);
    \end{tikzpicture}}
    \column{0.42\textwidth}
    \begin{block}{Seguridad}
      \small
      \begin{itemize}
        \item El ESP32 entra con un certificado X.509, sin contraseña
        \item La conexión la abre el ESP32, sin IP pública
        \item El PIN vive solo en la Lambda
        \item API limitada a 5 solicitudes por segundo
      \end{itemize}
    \end{block}
  \end{columns}
  \vspace{3mm}
  {\small Tópicos: \texttt{casa/puerta/eventos}, \texttt{casa/puerta/estado} y \texttt{casa/puerta/comandos}}
\end{frame}

\begin{frame}{Etapas 4 y 5 · Recorrido de un mensaje y métricas}
  \centering
  \resizebox{!}{0.62\textheight}{% Recorrido de un evento hasta la app, de una orden de abrir hasta el ESP32,
% y dónde se mide la latencia.
\begin{tikzpicture}[font=\footnotesize,
  cabeza/.style={caja=#1, minimum width=2.1cm, minimum height=0.8cm},
  vida/.style={draw=black!25, dashed},
  msg/.style={flecha},
  txt/.style={font=\scriptsize, text=black!75, fill=white, inner sep=1pt},
  fase/.style={font=\scriptsize\itshape, text=black!55, anchor=west},
]
  \node[cabeza=cProc] (e) at (0,0)    {ESP32};
  \node[cabeza=cApp]  (i) at (2.9,0)  {IoT Core};
  \node[cabeza=cApp]  (d) at (5.8,0)  {DynamoDB};
  \node[cabeza=cApp]  (l) at (8.7,0)  {API Gateway\\+ Lambda};
  \node[cabeza=cApp]  (w) at (11.6,0) {App móvil};
  \foreach \n in {e,i,d,l,w} \draw[vida] (\n.south) -- ++(0,-5.3);

  % Un evento sube y la app lo consulta
  \node[fase] at (-1.1,-0.85) {Evento};
  \draw[msg] (0,-1.2) -- node[txt, above] {evento (TLS, 8883)} (2.9,-1.2);
  \draw[msg] (2.9,-1.7) -- node[txt, above] {regla: guarda} (5.8,-1.7);
  \draw[msg, dashed] (11.6,-2.3) -- node[txt, above] {GET /eventos} (8.7,-2.3);
  \draw[msg] (8.7,-2.8) -- node[txt, above] {consulta} (5.8,-2.8);
  \draw[msg] (8.7,-3.3) -- node[txt, above] {eventos} (11.6,-3.3);

  % La app abre la puerta
  \draw[black!15] (-1.1,-3.7) -- (12.7,-3.7);
  \node[fase] at (-1.1,-4.0) {Abrir};
  \draw[msg, dashed] (11.6,-4.3) -- node[txt, above] {POST /abrir + PIN} (8.7,-4.3);
  \draw[msg] (8.7,-4.8) -- node[txt, above, pos=0.3] {orden} (2.9,-4.8);
  \draw[msg] (2.9,-5.3) -- node[txt, above] {servo 3\,s} (0,-5.3);

  % Marcas de tiempo
  \fill[cProc] (0,-1.2) circle (2pt) node[left, font=\scriptsize, text=cProc] {$t_e$};
  \fill[cApp] (2.9,-1.7) circle (2pt) node[left, font=\scriptsize, text=cApp, yshift=-4pt] {$t_r$};
  \node[font=\scriptsize, align=left, anchor=north west] at (-1.1,-5.6)
    {Latencia $= t_r - t_e$\quad($t_e$ con hora NTP, $t_r$ la agrega la regla)};
\end{tikzpicture}
}
  \vspace{2mm}
  \begin{columns}[T]
    \column{0.45\textwidth}
    \small
    \begin{itemize}
      \item Latencia, pérdida, disponibilidad
      \item Eventos por sensor, falsas alarmas, RSSI
    \end{itemize}
    \column{0.5\textwidth}
    \small
    \begin{itemize}
      \item El cálculo automático de las métricas queda para una siguiente etapa
      \item La app consulta cada 10 a 30\,s
    \end{itemize}
  \end{columns}
\end{frame}

\begin{frame}{Aplicación móvil}
  \centering
  \imagen{app.png}{0.85\linewidth}{0.75\textheight}{Capturas de la aplicación móvil puestas lado a lado: estado del nodo, lista de eventos y teclado del PIN.}
\end{frame}

\begin{frame}{Armado del prototipo}
  \begin{columns}[T]
    \column{0.32\textwidth}\fotopropia{prototipo.jpeg}{2.45cm}{Casa terminada}{Cartón reciclado}{Casa terminada, vista frontal}
    \column{0.32\textwidth}\fotopropia{compartimento.jpeg}{2.45cm}{Electrónica en la base}{Oculta bajo la casa}{Base abierta con ESP32, protoboard y cableado}
    \column{0.32\textwidth}\fotopropia{puerta.jpeg}{2.45cm}{Puerta}{Reed e imán}{Puerta abierta con el reed y el imán}
  \end{columns}
  \vspace{2mm}
  \begin{columns}[T]
    \column{0.32\textwidth}\fotopropia{sw420.jpeg}{2.45cm}{SW-420}{Pared izquierda}{SW-420 pegado por dentro}
    \column{0.32\textwidth}\fotopropia{servo.jpeg}{2.45cm}{Servo y pestillo}{En el piso de la casa}{Servo con el pestillo trabando la puerta}
    \column{0.32\textwidth}\fotopropia{prueba.jpg}{2.45cm}{Pruebas}{Equipo trabajando}{Equipo probando un evento crítico con la app abierta}
  \end{columns}
\end{frame}

% ===============================================================
\section{Resultados}

\begin{frame}{Pruebas del backend en AWS}
  \small
  \renewcommand{\arraystretch}{1.25}
  \centering
  \begin{tabular}{@{}cl c@{}}
    \toprule
    \textbf{N.°} & \textbf{Prueba, sin el ESP32} & \textbf{Resultado} \\
    \midrule
    1 & Evento publicado a mano llega a \texttt{EventosPuerta} & \textcolor{cAct}{\textbf{Correcto}} \\
    2 & PIN correcto, respuesta 200 y orden en \texttt{casa/puerta/comandos} & \textcolor{cAct}{\textbf{Correcto}} \\
    3 & PIN incorrecto, respuesta 401 y registro \texttt{pin\_incorrecto} & \textcolor{cAct}{\textbf{Correcto}} \\
    4 & \texttt{consultarPuerta} responde estado y eventos & \textcolor{cAct}{\textbf{Correcto}} \\
    5 & \texttt{GET /estado} y \texttt{GET /eventos} desde el navegador & \textcolor{cAct}{\textbf{Correcto}} \\
    6 & \texttt{POST /abrir} con \texttt{curl} & \textcolor{cAct}{\textbf{Correcto}} \\
    \bottomrule
  \end{tabular}
\end{frame}

\begin{frame}{Escenarios funcionales}
  \small
  \renewcommand{\arraystretch}{1.1}
  \centering
  \begin{tabular}{@{}cl>{\raggedright\arraybackslash}p{5.2cm}@{}}
    \toprule
    \textbf{N.°} & \textbf{Escenario} & \textbf{Comportamiento del programa} \\
    \midrule
    1 & Solo PIR & Alerta leve \\
    2 & Golpe y apertura en menos de 10\,s & Buzzer y pestillo \\
    3 & PIR y apertura con más de 10\,s & Dos alertas leves \\
    4 & Abrir desde la app con PIN y cerrar & Pestillo suelto y se traba al cerrar \\
    5 & Puerta abierta más de 1\,min & Buzzer hasta cerrarla \\
    6 & Crítico sin Wi-Fi & Alarma local igual \\
    7 & Servo y buzzer a la vez & El ESP32 no se reinicia \\
    8 & Humo frente al MQ-2 & Buzzer y pestillo suelto \\
    9 & Abrir sin orden previa & Evento \texttt{puerta\_forzada} \\
    \bottomrule
  \end{tabular}
\end{frame}

\begin{frame}{Sesión de prueba}
  \begin{columns}[c]
    \column{0.52\textwidth}
    \centering
    \imagen{prueba.jpg}{\linewidth}{5.5cm}{Equipo probando un evento crítico en el prototipo, con la aplicación abierta en un celular.}
    \column{0.46\textwidth}
    \begin{block}{Historial del 13 de noviembre}
      \small
      \begin{itemize}
        \item 22:38:19, movimiento, leve
        \item 22:40:42, PIN incorrecto
        \item 22:41:05, apertura desde la app
        \item 22:44:14, vibración, leve
        \item 22:44:18, puerta forzada, crítica
      \end{itemize}
    \end{block}
    {\small Golpe y apertura con 4\,s de diferencia, dentro de la ventana de 10\,s.}
  \end{columns}
\end{frame}

% ===============================================================
\section{Mercado}

\begin{frame}{Mercado y modelo de negocio}
  \centering
  \resizebox{!}{1.9cm}{% Modelo de negocio propuesto, en tres pasos.
\begin{tikzpicture}[font=\footnotesize,
  bloque/.style={caja=#1, minimum width=3.3cm, minimum height=1.5cm, text width=3.1cm},
  titulo/.style={font=\scriptsize\bfseries, text=#1},
]
  \node[bloque=cSensor] (a) at (0,0)   {Hogares sin alarma\\{\scriptsize casas, bodegas, cuartos alquilados}};
  \node[bloque=cAct]    (b) at (4.4,0) {Kit instalado\\{\scriptsize sensores, alarma y pestillo}};
  \node[bloque=cApp]    (c) at (8.8,0) {Suscripción\\{\scriptsize app con alertas y métricas}};
  \draw[flecha] (a) -- (b);
  \draw[flecha] (b) -- (c);
  \node[titulo=cSensor, above=2pt of a] {¿Para quién?};
  \node[titulo=cAct,    above=2pt of b] {Pago único};
  \node[titulo=cApp,    above=2pt of c] {Pago mensual};
\end{tikzpicture}
}
  \vspace{2mm}
  \begin{columns}[c]
    \column{0.5\textwidth}
    \centering
    \resizebox{!}{4.2cm}{% Costo por bloque (S/), con los precios de la lista de materiales (16 oct. 2026).
% Procesamiento: ESP32 40. Adquisición: PIR 8 + MC-38 5 + SW-420 6 + MQ-2 10 = 29.
% Conexiones: protoboard 7 + cables Dupont 8 + resistencias 2 + LED 0,6 = 17,6.
% Actuación (precios estimados, no están en la lista): buzzer 3 + servo SG90 12 = 15.
% Alimentación: cable USB 6.
\begin{tikzpicture}
\begin{axis}[grafico, xbar, bar width=0.42cm,
  width=\anchografico, height=\altografico,
  xmajorgrids, ymajorgrids=false,
  symbolic y coords={ali,act,con,adq,pro},
  ytick=data, yticklabels={Alimentación, Actuación (est.), Conexiones, Adquisición, Procesamiento}, xmin=0, xmax=50,
  xlabel={Costo (S/)},
  nodes near coords, every node near coord/.append style={font=\scriptsize, text=tinta},
  enlarge y limits=0.15]
  \addplot[fill=serie1, draw=white] coordinates
    {(6,ali) (15,act) (18,con) (29,adq) (40,pro)};
\end{axis}
\end{tikzpicture}
}
    \column{0.45\textwidth}
    \small
    \begin{itemize}
      \item Electrónica del nodo por unos S/~108, con buzzer y servo estimados
      \item La alarma local funciona sin suscripción
      \item La app muestra lo que una alarma barata no muestra
      \item La nube cobra por uso, sin un servidor encendido
    \end{itemize}
  \end{columns}
\end{frame}

\ifx\videourl\empty
\begin{frame}{Video publicitario}
  \centering
  \imagen{video.jpg}{0.75\linewidth}{6cm}{Fotograma del video con los cuatro integrantes y el prototipo.}
\end{frame}
\else
\begin{frame}{Video publicitario}
  \begin{columns}[c]
    \column{0.6\textwidth}
    \centering
    \imagen{video.jpg}{\linewidth}{5cm}{Fotograma del video con los cuatro integrantes y el prototipo.}
    \column{0.36\textwidth}
    \centering
    \qrcode[height=3.2cm]{\videourl}\par\vspace{2mm}
    {\footnotesize\url{\videourl}}
  \end{columns}
\end{frame}
\fi

\begin{frame}{Demostración en vivo}
  \begin{columns}[c]
    \column{0.6\textwidth}
    \begin{tikzpicture}[font=\small, node distance=5mm]
      \node[paso] (p1) {1};
      \node[right=3mm of p1, text width=6.6cm] {Pasar la mano frente al PIR. En la app aparece una \textbf{alerta leve}.};
      \node[paso, below=8mm of p1] (p2) {2};
      \node[right=3mm of p2, text width=6.6cm] {Golpear la puerta y abrirla. Suena el \textbf{buzzer} y el \textbf{pestillo} se traba.};
      \node[paso, below=8mm of p2] (p3) {3};
      \node[right=3mm of p3, text width=6.6cm] {Apagar el Wi-Fi y repetir. La alarma local sigue funcionando.};
      \node[paso, below=8mm of p3] (p4) {4};
      \node[right=3mm of p4, text width=6.6cm] {Abrir desde la app con el PIN. El pestillo se suelta, y al abrir y cerrar la puerta se traba solo.};
      \node[paso, below=8mm of p4] (p5) {5};
      \node[right=3mm of p5, text width=6.6cm] {Acercar humo al MQ-2. Suena el buzzer y la puerta se destraba.};
    \end{tikzpicture}
    \column{0.36\textwidth}
    \begin{alertblock}{Red para la demo}
      \small El Wi-Fi de la universidad pide iniciar sesión en una página, y el ESP32 no puede hacerlo. Usamos el punto de acceso de un celular.
    \end{alertblock}
  \end{columns}
\end{frame}

% ===============================================================
\section{Conclusiones}

\begin{frame}{Conclusiones y recomendaciones}
  \begin{columns}[T]
    \column{0.5\textwidth}
    \begin{block}{Conclusiones}
      \small
      \begin{itemize}
        \item La regla de 10\,s separa alertas leves de eventos críticos con una lógica simple
        \item La alarma y el bloqueo no dependen de Internet
        \item El peor caso bajó de unos 27\,W a unos 7\,W, con todo a 5\,V por USB
        \item AWS IoT Core y Lambda quitaron el servidor que había que mantener
        \item Falta medir la latencia p95 frente a la meta de 2\,s
      \end{itemize}
    \end{block}
    \column{0.46\textwidth}
    \begin{block}{Recomendaciones}
      \small
      \begin{itemize}
        \item Batería de respaldo
        \item Inicio de sesión en la API con Amazon Cognito
        \item Métricas, falsas alarmas y armar o silenciar en la app
        \item PCB propia en lugar del protoboard
        \item Varios nodos, ESP32-CAM y detección por CSI \cite{zhuang2021intrusion}
      \end{itemize}
    \end{block}
  \end{columns}
\end{frame}

% ===============================================================
\begin{frame}{Referencias}
  \tiny
  \begin{thebibliography}{10}
  \bibitem{alfuqaha2015survey} A. Al-Fuqaha et al., ``Internet of Things: A survey on enabling technologies, protocols, and applications,'' \emph{IEEE Commun. Surveys Tuts.}, vol. 17, no. 4, pp. 2347--2376, 2015.
  \bibitem{oktafiani2025esp32} F. Oktafiani y D. Widhiantoro, ``Studi literatur penggunaan ESP32 untuk sistem keamanan lingkungan rumah,'' \emph{SNIV}, vol. 4, no. 1, pp. 267--274, 2025.
  \bibitem{arinde2024multisensor} V. Arinde y L. Idowu, ``Multi-sensor intrusion detection system,'' arXiv:2406.05137, 2024.
  \bibitem{pugh2015fusion} M. Pugh, J. Brewer y J. Kvam, ``Sensor fusion for intrusion detection under false alarm constraints,'' \emph{IEEE SAS}, pp. 1--6, 2015.
  \bibitem{ray2018survey} P. P. Ray, ``A survey on Internet of Things architectures,'' \emph{J. King Saud Univ. Comput. Inf. Sci.}, vol. 30, no. 3, pp. 291--319, 2018.
  \bibitem{nannung2025iot} J. Nannung et al., ``Development of an IoT-based home security system prototype using multisensors and ESP32-CAM,'' \emph{JESSI}, vol. 6, no. 2, pp. 168--177, 2025.
  \bibitem{light2017mosquitto} R. A. Light, ``Mosquitto: server and client implementation of the MQTT protocol,'' \emph{JOSS}, vol. 2, no. 13, p. 265, 2017.
  \bibitem{aws2026iot} Amazon Web Services, ``What is AWS IoT?,'' \emph{AWS IoT Core Developer Guide}, consultado en nov. 2026.
  \bibitem{shi2016edge} W. Shi et al., ``Edge computing: Vision and challenges,'' \emph{IEEE Internet Things J.}, vol. 3, no. 5, pp. 637--646, 2016.
  \bibitem{aprilia2025boarding} M. W. Aprilia y A. Purwanto, ``IoT-based boarding room security system using PIR sensor and ESP32-CAM with Telegram notification,'' \emph{J. Ilmu Fisika dan Terapannya}, vol. 12, no. 2, pp. 95--102, 2025.
  \bibitem{jama2026edge} Y. A. Jama, K. W. Muse y R. M. Ali, ``Empirical evaluation of a layered edge--cloud IoT security architecture for resilient residential access control,'' \emph{Open Res. Africa}, vol. 9, p. 35, 2026.
  \bibitem{mishra2020mqtt} B. Mishra y A. Kertesz, ``The use of MQTT in M2M and IoT systems: A survey,'' \emph{IEEE Access}, vol. 8, pp. 201071--201086, 2020.
  \bibitem{bender2021mqtt} M. Bender et al., ``Open-source MQTT evaluation,'' \emph{IEEE CCNC}, pp. 1--4, 2021.
  \bibitem{zhuang2021intrusion} W. Zhuang et al., ``Develop an adaptive real-time indoor intrusion detection system based on empirical analysis of OFDM subcarriers,'' \emph{Sensors}, vol. 21, no. 7, art. 2287, 2021.
  \end{thebibliography}
\end{frame}

\end{document}
